% praeambel.tex — Paket- und Layout-Konfiguration
% Whitepaper „Agentic Software Development — Enterprise Architecture with
% AI Agents", modernisierte Fassung 2026 (v1.3 -> v2.x/v3.x)

% --- Encoding, Sprache, Schrift ---
\usepackage[T1]{fontenc}
\usepackage{lmodern}
% Copy-Paste-Fidelity: ToUnicode-CMaps fuer alle Glyphen. Ohne dies liefert
% das Kopieren aus Code-Listings u. a. fuer ^ den Modifier-Zirkumflex ˆ
% (U+02C6) statt U+005E — eine kopierte Regex-Zeichenklasse [ˆ"] waere
% dann nicht mehr invertierend.
\input{glyphtounicode}
\pdfgentounicode=1
\usepackage[ngerman]{babel}
\usepackage[style=german]{csquotes}
\usepackage{microtype}

% --- Zeilenabstand ---
\usepackage{setspace}
\onehalfspacing

% --- Seitenlayout ---
\usepackage[a4paper,margin=2.5cm]{geometry}

% --- Mathematik ---
\usepackage{amsmath}

% --- Tabellen ---
\usepackage{booktabs}
% Pandoc setzt Markdown-Pipe-Tabellen als longtable (braucht zusätzlich array/calc)
\usepackage{longtable}
\usepackage{array}
\usepackage{calc}

% --- Grafiken ---
\usepackage{graphicx}
\graphicspath{{abbildungen/out/}}
% Pandoc emittiert width UND height; ohne keepaspectratio würden die
% Abbildungen auf \textheight verzerrt (Pandocs eigene Preamble setzt dies auch).
\setkeys{Gin}{keepaspectratio}

% --- Bibliographie ---
\usepackage[
  backend=biber,
  style=numeric,
  sorting=nyt,
  urldate=long,
  dateabbrev=false
]{biblatex}
\addbibresource{literatur.bib}

% --- Hyperlinks (dezent, keine farbigen Rahmen) ---
\usepackage[hidelinks]{hyperref}
% xurl erlaubt Zeilenumbrüche an beliebigen Stellen in URLs (Literaturverzeichnis)
\usepackage{xurl}

% --- Codeblöcke: lange Zeilen umbrechen statt in den Rand zu laufen ---
% Pandoc setzt Code als verbatim; fvextra rüstet Zeilenumbruch nach.
% breakautoindent: die Fortsetzungszeile beginnt in der Einrückspalte der
% umbrochenen Zeile; ein dezenter grauer Pfeil markiert den Umbruch.
% upquote: gerade Anführungszeichen und Backticks in Verbatim (statt
% typographischer) — kopierter Code bleibt gültiger Code.
\usepackage{fvextra}
\usepackage{upquote}
\usepackage{xcolor}
% accsupp: der Umbruchpfeil ist reine Optik und darf beim Kopieren aus dem
% PDF keinen Text beitragen (ActualText leer).
\usepackage{accsupp}
\newcommand{\umbruchpfeil}{%
  \BeginAccSupp{ActualText={}}%
  \textcolor{gray}{\tiny$\hookrightarrow$}%
  \EndAccSupp{}}
\RecustomVerbatimEnvironment{verbatim}{Verbatim}{%
  breaklines=true,
  breakautoindent=true,
  breaksymbolleft={\umbruchpfeil},
  breaksymbolindentleft=0.6em,
  fontsize=\small}

% Pandoc-Syntax-Highlighting (Stil: tango; Makros generiert via
% `pandoc -s --highlight-style=tango`, abgelegt in praeambel-highlighting.tex)
\input{praeambel-highlighting}
\RecustomVerbatimEnvironment{Highlighting}{Verbatim}{%
  commandchars=\\\{\},
  breaklines=true,
  breakautoindent=true,
  breaksymbolleft={\umbruchpfeil},
  breaksymbolindentleft=0.6em,
  fontsize=\small}

% --- Praxis-Check-Kaesten und Hinweis-Bloecke zusammenhalten ---
% Pandoc setzt Blockquotes als quote-Umgebung. Ohne Mindestplatz-Anforderung
% reisst ein kurzer Kasten am Seitenende auf (eine Zeile auf der Folgeseite).
\usepackage{needspace}
\usepackage{etoolbox}
\AtBeginEnvironment{quote}{\Needspace*{4\baselineskip}}

% --- Float-Platzierung: Abbildungen naeher am Text halten ---
% Voreinstellung von LaTeX schiebt grosse Abbildungen leicht mehrere Seiten
% weit; in einem abbildungsreichen Dokument reisst das Bild und Bezugstext
% auseinander.
\renewcommand{\topfraction}{0.9}
\renewcommand{\bottomfraction}{0.8}
\renewcommand{\textfraction}{0.07}
\renewcommand{\floatpagefraction}{0.7}
\setcounter{topnumber}{3}
\setcounter{bottomnumber}{2}
\setcounter{totalnumber}{4}
% placeins stellt \FloatBarrier bereit — punktuell dort eingesetzt, wo eine
% verschobene Abbildung sonst mitten in einer mehrseitigen Tabelle landet
% (z. B. Komponenten-Tabelle in Kapitel 12).
\usepackage{placeins}
% Mehrseitige Tabellen (longtable) nicht mit Kopf + einzelner Zeile am
% Seitenende beginnen lassen — dann lieber die Tabelle auf die Folgeseite.
\BeforeBeginEnvironment{longtable}{\Needspace*{6\baselineskip}}

% --- Inhaltsverzeichnis: Platz fuer zweistellige Unterabschnittsnummern ---
% Ohne Verbreiterung klebt "19.10" im Sonderdruck am Titel ("19.10Vom Run").
\DeclareTOCStyleEntry[numwidth=3.1em]{tocline}{section}
\DeclareTOCStyleEntry[numwidth=2.0em]{tocline}{chapter}
% Gleiches Problem bei den Teilen: die Voreinstellung ist fuer dreistellige
% roemische Ziffern zu schmal ("VII.Architekturentscheidungen").
\DeclareTOCStyleEntry[numwidth=3.0em]{tocline}{part}
% Gleiches Problem im Abbildungs- und Tabellenverzeichnis: dort tragen die
% Eintraege Kapitelnummer + laufende Nummer ("19.10."), was in der schmalen
% Voreinstellung ohne Abstand an den Titel stiess ("19.10.Deployment-Sicht").
\DeclareTOCStyleEntry[numwidth=3.1em]{tocline}{figure}
\DeclareTOCStyleEntry[numwidth=3.1em]{tocline}{table}

% --- Pandoc-Kompatibilität ---
% \tightlist wird von Pandoc für kompakte Listen ohne Zusatzabstand erzeugt
% und ist ohne den Pandoc-Standardtemplate nicht definiert.
\providecommand{\tightlist}{%
  \setlength{\itemsep}{0pt}\setlength{\parskip}{0pt}}

% Einzelne Unicode-Zeichen aus wörtlich zitierten Roh-Extrakten, die im
% Standard-T1-Encoding kein Glyphen-Mapping haben.
\DeclareUnicodeCharacter{03A3}{$\Sigma$}
\DeclareUnicodeCharacter{2217}{$\ast$}
% Zeichen aus dem Änderungsvermerk (Anhang B)
\DeclareUnicodeCharacter{2192}{$\rightarrow$}
\DeclareUnicodeCharacter{2194}{$\leftrightarrow$}
\DeclareUnicodeCharacter{2265}{$\geq$}
% Zeichen aus der PDF-Extraktion des Agentic-Software-Development-Whitepapers
% (Tabellen, Formeln in ADR-1 „n×(n-1)/2", Verzeichnisbaum in ADR-4)
\DeclareUnicodeCharacter{00D7}{$\times$}
\DeclareUnicodeCharacter{2248}{$\approx$}
\DeclareUnicodeCharacter{2500}{-}   % ─ Box-Drawing: horizontale Linie
\DeclareUnicodeCharacter{2502}{|}   % │ Box-Drawing: vertikale Linie
\DeclareUnicodeCharacter{2514}{+}   % └ Box-Drawing: Ecke unten links
\DeclareUnicodeCharacter{251C}{+}   % ├ Box-Drawing: T-Stück links

% --- Abkürzungsverzeichnis (Setup vorbereitet, Liste vorerst leer) ---
\usepackage[printonlyused]{acronym}
% Beispiel-Einträge (auskommentiert, bei Bedarf befüllen — Kap. 23 „Glossar
% & Abkürzungsverzeichnis" führt die Abkürzungen bereits als Markdown-Tabelle):
% \begin{acronym}
%   \acro{DDD}{Domain-Driven Design}
%   \acro{ADR}{Architecture Decision Record}
% \end{acronym}
